\begin{textsample}{POS Dim 1 – pe – Score 46.00 – pe/pe\_ef\_47.txt}  \label{ex:f1_pos_241}
4.2.2 A MATEMÁTICA NA SALA DE \textbf{AULA} Ao considerar os processos de ensino e aprendizagem que ocorrem no interior da sala de \textbf{aula}, parte -se do pressuposto de que aprender \textbf{implica} na construção dos conceitos pelo próprio estudante, na medida em que ele é desafiado a confrontar antigas concepções, inclusive da sua vivência extraescolar, e levado a elaborar novos conceitos \textbf{esperados} pela escola.
Essa concepção apresenta uma lógica diferente da comum, ou seja, a aprendizagem de um novo conceito acontece pela apresentação de uma situação-problema ao estudante, \textbf{instigando} o mesmo à \textbf{compreensão} \textbf{conceitual}.
A \textbf{análise} dessa situação conduz à definição, à generalização e à \textbf{sistematização} do conceito, que vai sendo construído ao longo do processo de aprendizagem.
Por sua vez, os mesmos conceitos são retomados, posteriormente, e aprofundados em níveis mais \textbf{complexos}, de forma a conduzir o estudante a relacionar o que já sabia com o que virá a aprender em um novo contexto ( CÂMARA e LIMA, 2010 ).
É importante ressaltar que uma situação sem sentido não pode levar a uma aprendizagem consistente e duradoura.
Recomenda -se propor e explorar atividades matemáticas ricas e produtivas, considerando as experiências e os interesses dos estudantes.
Com base nos \textbf{Parâmetros} Curriculares de Matemática ( PERNAMBUCO, 2012 ), deve -se tomar como ponto de partida a ideia de que aprender Matemática vai além de simplesmente acumular conteúdos.
O estudante deve ser conduzido a"fazer"Matemática.
Cada vez mais, defende -se a ideia de que é preciso saber e saber fazer Matemática.
No contexto dessa discussão, mesmo que de forma simplificada, deve -se associar o saber aos conhecimentos apreendidos pelo estudante, e o saber fazer à sua \textbf{capacidade} de \textbf{mobilizar} esses conhecimentos como resposta a um \textbf{problema} ( CÂMARA e LIMA, 2010 ).
A ideia de fazer matemática \textbf{exige} esforço, engajamento e iniciativa.
A sala de \textbf{aula} deve ser um ambiente onde os estudantes sejam convidados a buscar soluções para os \textbf{problemas} apresentados, conduzindo -os a pensar, \textbf{argumentar} e dar sentido.
É importante criar um espaço no qual os estudantes devem ser \textbf{instigados} a compreender ativamente os conceitos matemáticos explorados, testar ideias e fazer conjecturas, desenvolver \textbf{raciocínios} e apresentar explicações de forma escrita ou verbal.
Para isso, são oferecidos diferentes \textbf{caminhos} ao professor, tais como a estratégia da resolução de \textbf{problema}, a investigação, a modelagem matemática, as \textbf{tecnologias} \textbf{digitais}, a calculadora, a evolução histórica dos conceitos matemáticos, os jogos matemáticos na sala de \textbf{aula}, o desenvolvimento de projetos de trabalho colaborativo, a etnomatemática ou \textbf{abordagem} cultural, entre outros.
Essas formas privilegiadas da atividade matemática são, ao mesmo tempo, objeto e estratégia para a aprendizagem ao longo de toda a educação básica.
Esses processos de aprendizagem são potencialmente ricos para o desenvolvimento de \textbf{competências} fundamentais, por exemplo, para o \textbf{letramento} matemático ( \textbf{raciocínio}, representação, comunicação e argumentação ) e para o desenvolvimento do pensamento computacional.
A literatura tem constatado que, para grande parte dos conceitos e procedimentos trabalhados na escola, a aprendizagem não se realiza em um único período nem em um espaço muito limitado de tempo.
É importante levar em conta que a aprendizagem é mais eficiente quando os conteúdos são revisitados, de forma progressivamente ampliada e aprofundada, durante todo o percurso escolar.
É \textbf{pertinente} destacar ainda a ideia de que os conceitos \textbf{relevantes} para a formação matemática \textbf{atual} deve ser abordada desde o início da formação escolar.
Tal ponto de \textbf{vista} apoia -se na concepção de que a construção de um conceito pelas pessoas é um processo dinâmico que se completa no \textbf{decorrer} de um longo período, de estágios mais intuitivos aos mais formais.
O ensino da Matemática deve estar em sintonia com as práticas matemáticas científicas e sociais fora da escola, bem como se apoiar de forma adequada nos meios tecnológicos que instrumentalizam essas práticas.
Ao mesmo tempo, é fundamental reconhecer que a construção desses conhecimentos não acontece de \textbf{maneira} espontânea, mas como consequência da \textbf{mobilização} de recursos \textbf{metodológicos} adequados e estratégias de ensino produtivas, que serão abordados e apresentados posteriormente.
Nos dias \textbf{atuais}, também é fundamental considerar para a aprendizagem as utilizações reais e as potencialidades oferecidas pelas \textbf{tecnologias} \textbf{digitais}.
É \textbf{pertinente} também destacar que a resolução de \textbf{problemas} seja considerada o \textbf{foco} principal a nortear o ensino da matemática na etapa do ensino fundamental \textbf{tanto} nos anos iniciais como nos anos \textbf{finais}.
Tal estratégia de ensino permite que as atividades ou \textbf{problemas} podem e devem ser propostos de modo a envolver os estudantes no pensar, no fazer e no desenvolver a matemática básica e essencial tão necessária para a sua aprendizagem.
Para isto, é importante planejar e propor atividades motivadoras que \textbf{instiguem} a curiosidade dos estudantes, que os levem a investigar, a experimentar, a confrontar e interpretar resultados, de modo a buscar respostas ou soluções para as situações vivenciadas dentro e fora da sala de \textbf{aula}.
As atividades mais eficazes partem de onde os estudantes estão, dos seus conhecimentos prévios.
Isto é, para ensinar, deve -se iniciar com as ideias que as crianças já possuem – as que servirão de ponto de partida para criar novas ideias.
Portanto, envolver e engajar os estudantes requer \textbf{tarefas} ou atividades que sejam fundamentadas em \textbf{problemas}, pois eles aprendem matemática como resultado da experiência de resolução de \textbf{problemas} em vez de elementos que devem ser ensinados antes de resolver \textbf{problemas}.
Sendo assim, o processo de resolução de \textbf{problemas} está articulado com a aprendizagem, cujo efeito \textbf{esperado} é o seguinte : as crianças aprendem matemática fazendo matemática.
Na perspectiva didática, pensa -se no trabalho de modelagem na sala de \textbf{aula} como \textbf{caminho} para que os estudantes tenham uma experiência de produção do conhecimento no campo de \textbf{certo} domínio matemático, experiência que possa permitir, também, enriquecer a conceitualização \textbf{teórica} nesse mesmo domínio.
Isso \textbf{demanda} que seja analisado cada domínio que é objeto de ensino, levando em consideração os \textbf{problemas} que os conceitos permitem abordar, as propriedades que relacionam os conceitos e que se traduzem em estratégias de resolução na medida em que possibilitam transformar as relações envolvidas em um \textbf{problema} e as formas de representação que têm destaque ( SADOVSKY, 2007 ).
Na cultura escolar, nas últimas décadas, os conteúdos matemáticos a serem ensinados e aprendidos têm sido organizados em grandes campos.
Entretanto, é \textbf{indispensável} que tais campos não sejam vistos como blocos estanques e autossuficientes.
Denominamos esses campos de conhecimentos no componente curricular de Matemática por unidades temáticas.
Neste documento, para as duas etapas do ensino fundamental - anos iniciais e anos \textbf{finais} - as habilidades relativas às aprendizagens esperadas serão apresentadas pelas seguintes unidades temáticas : Geometria, Estatística e Probabilidade, Álgebra, Grandezas e Medidas e Números.
No entanto, é preciso ter clareza de que essa organização se deve, unicamente, à busca de uma melhor forma de apresentação, pois, no trabalho em sala de \textbf{aula}, é importante que o professor busque, de forma sistemática, articular as unidades temáticas, trabalhando de forma integrada, considerando os objetos de conhecimento e habilidades previstos por ano de escolaridade.
Na unidade temática Geometria, o estudo de posição e deslocamentos no espaço, das \textbf{figuras} geométricas e das relações entre elementos de \textbf{figuras} planas e espaciais contribui para o desenvolvimento do pensamento geométrico por parte dos estudantes.
Esse pensamento é necessário para investigar propriedades, fazer conjecturas e produzir \textbf{argumentos} geométricos convincentes, ao mesmo tempo que compreende um conjunto de conceitos e procedimentos para resolver \textbf{problemas} do mundo físico e de diferentes áreas do conhecimento.
Destacam -se as ideias matemáticas fundamentais associadas a essa temática que são, principalmente, construção, representação e interdependência.
Em Probabilidade e Estatística são estudados a incerteza e o tratamento de dados/informações.
Essa unidade temática propõe a \textbf{abordagem} de conceitos, fatos e procedimentos presentes em muitas situações-problema da vida cotidiana, das \textbf{ciências} e da \textbf{tecnologia}.
Assim, todos os cidadãos precisam desenvolver habilidades para coletar, organizar, representar, interpretar e analisar dados em uma variedade de contextos de \textbf{maneira} a fazer julgamentos bem fundamentados e tomar as decisões adequadas.
Merece destaque o uso de \textbf{tecnologias} como calculadoras e planilhas eletrônicas, que podem ser utilizadas como recursos para avaliar, comparar e organizar conjunto de dados em gráficos, bem como para efetuar cálculos e analisar as medidas de tendência \textbf{central}.
Além disto, trabalhar com as noções que sustentam o conceito de probabilidade como aleatoriedade e chance são fundamentais nessa unidade.
A unidade temática Álgebra tem como \textbf{foco} o desenvolvimento de um tipo especial de pensamento – pensamento algébrico – que é essencial para utilizar modelos matemáticos na \textbf{compreensão}, representação e \textbf{análise} de relações quantitativas de grandezas e, também, de situações e estruturas matemáticas, fazendo uso de letras e outros símbolos.
Para esse desenvolvimento, é necessário que os estudantes identifiquem regularidades e padrões de \textbf{sequências} numéricas e não numéricas, estabeleçam leis matemáticas que expressem relações de interdependência entre grandezas em diferentes contextos, bem como criar, interpretar e transitar entre as diversas representações gráficas e simbólicas para resolver \textbf{problemas} por meio de equações e inequações com \textbf{compreensão} dos procedimentos utilizados.
As ideias matemáticas fundamentais vinculadas a essa unidade são : equivalência, variação, interdependência e proporcionalidade.
Como as medidas quantificam grandezas do mundo físico e são fundamentais para a \textbf{compreensão} da realidade, a unidade temática Grandezas e Medidas propõe o estudo das medidas e das relações entre elas, favorecendo a integração da Matemática a outras áreas de conhecimento, como \textbf{Ciências} ( densidade, grandezas e escalas do Sistema Solar, energia elétrica etc. ) ou Geografia ( coordenadas geográficas, densidade demográfica, escalas de mapas e guias etc. ).
Essa unidade temática contribui ainda para a consolidação e ampliação da noção de número, a aplicação de noções geométricas e a construção do pensamento algébrico, assim como as grandezas e medidas são fortes articuladores com as práticas sociais e profissionais.
Além do papel articulador dessa unidade temática, é fundamental que os estudantes compreendam a noção de grandeza enquanto atributo de um objeto ; identifiquem diferentes grandezas associadas a um mesmo objeto ; percebam a diferença entre uma \textbf{figura}, as grandezas a ela associadas e o número associado a medição dessa grandeza ; saibam utilizar instrumentos de medição e compreendam a diferença entre a medição prática e a \textbf{teórica} e entre precisão, \textbf{erro} e estimativa de medidas.
A finalidade da unidade temática Números é desenvolver o pensamento numérico, que \textbf{implica} o conhecimento de \textbf{maneiras} de quantificar atributos de objetos e de \textbf{julgar} e interpretar \textbf{argumentos} \textbf{baseados} em quantidades.
No processo da construção da noção de número, os estudantes precisam desenvolver, entre outras, as ideias de aproximação, de proporcionalidade e de equivalência e ordem, noções fundamentais da Matemática.
Para essa construção, é importante propor, por meio de situações significativas, sucessivas ampliações dos campos numéricos.
No estudo desses campos numéricos, devem ser enfatizados registros, usos, significados e operações.

% matched lemmas: abordagem, análise, argumentar, argumento, atual, aula, basear, caminho, capacidade, central, certo, ciência, competência, complexo, compreensão, conceitual, decorrer, demandar, digital, erro, esperar, exigir, figura, final, foco, implicar, indispensável, instigar, julgar, letramento, maneira, metodológico, mobilizar, mobilização, parâmetro, pertinente, problema, raciocínio, relevante, sequência, sistematização, tanto, tarefa, tecnologia, teórico, vista
\end{textsample}
